\section{Discussion}
\label{sec:discussion}

The three preserved-utility methods in Table~\ref{tab:cifar10-ci-baselines} partition
along request-time cost and guarantee rather than along utility. FORGE processes an
erasure in $0.31$~s with $O(1)$ per-client server storage and a
$(\varepsilon{=}2.0, \delta{=}10^{-5})$-DP guarantee; SIFU takes $2130$~s with
$O(R|\theta|)$ storage and an $\varepsilon{=}1.0$ influence bound; FedRecovery at
$\varepsilon{=}\infty$ takes $16.4$~s with $O(NR|\theta|)$ storage and no DP
guarantee. SIFU and FORGE represent two design points on the same utility surface:
SIFU accepts request-time cost in exchange for an unmodified training loop; FORGE
accepts a slightly more complex training protocol in exchange for constant-latency
erasure. Which is preferable depends on the deployment's cost structure.

The FedRecovery $\varepsilon$-sweep is a structural result rather than a
particular-$\varepsilon$ artefact: every finite $\varepsilon \in \{1,2,5,10\}$
collapses to chance accuracy under a clip norm compatible with trainable FCL on
ResNet-18. The sole operating point that preserves utility ($\varepsilon{=}\infty$)
emits no DP guarantee and leaks at $0.111$ on Yeom. FedRecovery therefore has no
feasible point in the tested regime that simultaneously preserves utility and limits
leakage.
